\documentclass[11pt]{article}
% Shared preamble for Bliss documentation.
% Input from a document after \documentclass, before \begin{document}.

\usepackage[a4paper,margin=1in]{geometry}
\usepackage[T1]{fontenc}
\usepackage{lmodern}
\usepackage{microtype}
\usepackage{xcolor}
\usepackage{booktabs}
\usepackage{longtable}
\usepackage{array}
\usepackage{colortbl}
\usepackage{enumitem}
\usepackage{titlesec}
\usepackage{fancyhdr}
\usepackage{bytefield}
\usepackage{listings}
\usepackage[most]{tcolorbox}
\usepackage[hidelinks]{hyperref}

% ---------------------------------------------------------------------------
% Palette
% ---------------------------------------------------------------------------
\definecolor{blissink}{HTML}{1C2333}   % near-black text
\definecolor{blissaccent}{HTML}{3B4CCA} % indigo
\definecolor{blissaccent2}{HTML}{0E7C7B}% teal
\definecolor{blissrule}{HTML}{C7CBD6}
\definecolor{blisssoft}{HTML}{F1F2F6}
\definecolor{blisscode}{HTML}{F6F8FA}
\definecolor{blisskw}{HTML}{9B1D64}
\definecolor{blisscmt}{HTML}{6A737D}

\color{blissink}

% ---------------------------------------------------------------------------
% Headings
% ---------------------------------------------------------------------------
\titleformat{\section}
  {\normalfont\Large\bfseries\color{blissaccent}}
  {\thesection}{0.6em}{}
\titleformat{\subsection}
  {\normalfont\large\bfseries\color{blissink}}
  {\thesubsection}{0.6em}{}
\titlespacing*{\section}{0pt}{2.2ex plus 1ex minus .2ex}{1.2ex plus .2ex}

% ---------------------------------------------------------------------------
% Headers / footers
% ---------------------------------------------------------------------------
\pagestyle{fancy}
\fancyhf{}
\renewcommand{\headrulewidth}{0.4pt}
\renewcommand{\footrulewidth}{0pt}
\fancyhead[L]{\small\color{blisscmt}\textsc{Bliss}}
\fancyhead[R]{\small\color{blisscmt}\leftmark}
\fancyfoot[C]{\small\color{blisscmt}\thepage}

% ---------------------------------------------------------------------------
% Tables
% ---------------------------------------------------------------------------
\renewcommand{\arraystretch}{1.25}
\arrayrulecolor{blissrule}
\newcolumntype{L}[1]{>{\raggedright\arraybackslash}p{#1}}
\newcolumntype{C}[1]{>{\centering\arraybackslash}p{#1}}

% ---------------------------------------------------------------------------
% Monospace helper
% ---------------------------------------------------------------------------
\newcommand{\code}[1]{\texttt{#1}}
\newcommand{\reg}[1]{\texttt{\textbf{#1}}}

% ---------------------------------------------------------------------------
% Listings — Bliss assembly
% ---------------------------------------------------------------------------
\lstdefinelanguage{blissasm}{
  morekeywords={hlt,nop,mov,ldi16,ldi32l,ldi32h,ldm8,ldm16,ldm32,str8,str16,str32,
    add,sub,mul,and,or,xor,not,sll,slr,sar,beq,bne,blt,bge,bltu,bgeu,bal,balr,
    ecall,eret,mfcr,mtcr,li,push,pop,call,ret,jmp,addi,cmp},
  morecomment=[l];,
  morestring=[b]",
  sensitive=false,
}
\lstdefinestyle{bliss}{
  language=blissasm,
  backgroundcolor=\color{blisscode},
  basicstyle=\ttfamily\small,
  keywordstyle=\color{blisskw}\bfseries,
  commentstyle=\color{blisscmt}\itshape,
  stringstyle=\color{blissaccent2},
  frame=single,
  framerule=0.4pt,
  rulecolor=\color{blissrule},
  framesep=6pt,
  xleftmargin=6pt,
  xrightmargin=6pt,
  showstringspaces=false,
  breaklines=true,
  columns=fullflexible,
  keepspaces=true,
}
\lstset{style=bliss}

% ---------------------------------------------------------------------------
% Callout box
% ---------------------------------------------------------------------------
\newtcolorbox{blissnote}[1][Note]{
  enhanced, breakable,
  colback=blisssoft, colframe=blissaccent,
  boxrule=0pt, leftrule=3pt,
  arc=1pt, left=8pt, right=8pt, top=6pt, bottom=6pt,
  fonttitle=\bfseries\color{blissaccent},
  title={#1},
}

% ---------------------------------------------------------------------------
% bytefield defaults
% ---------------------------------------------------------------------------
\bytefieldsetup{
  bitwidth=1.15em,
  endianness=big,
  boxformatting={\centering\small},
  bitformatting=\scriptsize\color{blisscmt},
}


\title{\vspace{-2em}\color{blissaccent}\textbf{\Huge Bliss}\\[0.4em]
  \color{blissink}\Large ISA \& System Reference}
\author{}
\date{Generated from \texttt{SPEC.md}}

\begin{document}
\maketitle
\thispagestyle{empty}

\begin{center}
\begin{minipage}{0.82\linewidth}
\small\itshape
A minimal 32-bit register machine, an assembler for it, and a small
terminal-only operating system --- each layer hand-rolled so that every byte is
understood. This document is the reference for the virtual machine's
architecture: registers, memory, instructions, the privilege mechanism, the
calling convention, the syscall ABI, and the BlissFS on-disk format.
\end{minipage}
\end{center}

\vspace{1em}
\tableofcontents
\vspace{1em}

% ===========================================================================
\section{Architecture at a glance}
% ===========================================================================

\begin{itemize}[leftmargin=1.4em]
  \item 32-bit, register-based virtual machine.
  \item 16 general-purpose 32-bit registers (\code{r0}--\code{r15}).
  \item Fixed-width 32-bit instructions; the opcode is always bits \code{[31:26]}.
  \item Load/store architecture: arithmetic operates on registers only; memory
        is touched exclusively through explicit loads and stores.
  \item Memory-mapped I/O --- there are no dedicated I/O instructions.
  \item Flat 32-bit address space.
  \item Comparison-in-branch: there is no flags register; branch instructions
        take two source registers and compare them directly.
\end{itemize}

\begin{blissnote}[Design lineage]
The ISA draws on RISC-V and early ARM: fixed-width instructions, a load/store
core, a generous register file, branch instructions that fold in the comparison,
and MMIO for every device.
\end{blissnote}

% ===========================================================================
\section{Registers}
% ===========================================================================

\subsection{General-purpose file}

All sixteen registers are 32 bits wide. Three carry an architected role:

\begin{center}
\begin{tabular}{@{}L{2.2cm} L{11cm}@{}}
\toprule
\textbf{Register} & \textbf{Role} \\
\midrule
\reg{r0}          & Assembler scratch / syscall number / return value. Reserved:
                    never assume it survives a pseudo-instruction or a call. \\
\reg{r1}--\reg{r12} & Ordinary general-purpose registers (see the calling
                    convention for how they split into argument, caller-saved
                    and callee-saved groups). \\
\reg{r13}         & Stack pointer (\code{SP}). \\
\reg{r14}         & Link register (\code{LR}). \\
\reg{r15}         & Program counter (\code{PC}). Not directly writable. \\
\bottomrule
\end{tabular}
\end{center}

\subsection{Control registers}

Four control registers sit alongside the general-purpose file and are accessed
only with \code{MFCR} / \code{MTCR}.

\begin{center}
\begin{tabular}{@{}L{2.0cm} C{1.4cm} L{9.6cm}@{}}
\toprule
\textbf{Name} & \textbf{Index} & \textbf{Description} \\
\midrule
\reg{tvec}  & 0 & Trap vector: address of the supervisor trap handler. \\
\reg{epc}   & 1 & Exception PC: return address saved on trap. \\
\reg{cause} & 2 & Trap cause code (\code{0} = syscall, \dots). \\
\reg{mode}  & 3 & Current privilege level (\code{0} = user, \code{1} = supervisor). \\
\bottomrule
\end{tabular}
\end{center}

% ===========================================================================
\section{Privilege modes}
% ===========================================================================

Bliss has two privilege levels: \textbf{user} (0) and \textbf{supervisor} (1).
The current level lives in the \reg{mode} control register. \emph{The CPU boots
in supervisor mode.}

\subsection{Trap flow --- \code{ECALL}}

When \code{ECALL} executes:
\begin{enumerate}[leftmargin=1.8em,itemsep=0.1em]
  \item \code{epc $\leftarrow$ PC} --- the address of the instruction following \code{ECALL}.
  \item \code{cause $\leftarrow$ 0} --- syscall.
  \item \code{mode $\leftarrow$ supervisor}.
  \item \code{PC $\leftarrow$ tvec}.
\end{enumerate}

\subsection{Return flow --- \code{ERET}}

When \code{ERET} executes:
\begin{enumerate}[leftmargin=1.8em,itemsep=0.1em]
  \item \code{mode $\leftarrow$ user}.
  \item \code{PC $\leftarrow$ epc}.
\end{enumerate}

\begin{blissnote}[No automatic save/restore]
Neither \code{ECALL} nor \code{ERET} saves or restores any general-purpose
register. Preserving state across a trap is entirely the trap handler's
responsibility.
\end{blissnote}

% ===========================================================================
\section{Memory map}
% ===========================================================================

\begin{center}
\begin{tabular}{@{}L{4.4cm} L{9.0cm}@{}}
\toprule
\textbf{Address range} & \textbf{Use} \\
\midrule
\code{0x0000\_0000}--\code{0x0000\_0FFF} & Kernel code and data \\
\code{0x0000\_1000}--\code{0x0000\_11FF} & Kernel scratch buffer (disk I/O, one sector) \\
\code{0x0000\_1200}--\code{0x0000\_1FFB} & Kernel stack (grows downward; \code{SP} starts at \code{0x1FFC}) \\
\code{0x0000\_2000}-- & User program \\
\addlinespace[0.3em]
\code{0xFFFF\_0000} & Serial TX --- write byte $\rightarrow$ stdout \\
\code{0xFFFF\_0004} & Serial RX --- read byte $\leftarrow$ stdin (blocks until ready) \\
\code{0xFFFF\_0010} & Disk sector register \\
\code{0xFFFF\_0014} & Disk buffer register \\
\code{0xFFFF\_0018} & Disk command (\code{0} = read sector) \\
\bottomrule
\end{tabular}
\end{center}

\begin{blissnote}[Warning]
Memory protection is not implemented. User code may read or write any address,
including MMIO registers and kernel memory. This is a deliberate simplification
for now.
\end{blissnote}

% ===========================================================================
\section{Instruction set}
% ===========================================================================

\subsection{Data movement}
\begin{center}
\begin{tabular}{@{}L{3.0cm} L{10.4cm}@{}}
\toprule
\code{MOV}            & Register to register. \\
\code{LDI16}          & Load 16-bit immediate (zero-extended). \\
\code{LDI32L}         & Load immediate into the lower 16 bits of a register. \\
\code{LDI32H}         & Load immediate into the upper 16 bits of a register. \\
\code{LDM8/16/32}     & Load from memory into a register (8/16 zero-extended). \\
\code{STR8/16/32}     & Store a register to memory. \\
\bottomrule
\end{tabular}
\end{center}

\subsection{Arithmetic and logic}
\begin{center}
\begin{tabular}{@{}L{3.0cm} L{10.4cm}@{}}
\toprule
\code{ADD}, \code{SUB}, \code{MUL} & Integer arithmetic. \\
\code{AND}, \code{OR}, \code{XOR}, \code{NOT} & Bitwise logic. \\
\code{SLL}, \code{SLR}            & Logical left / right shift. \\
\code{SAR}                        & Arithmetic right shift. \\
\bottomrule
\end{tabular}
\end{center}

\subsection{Control flow}
\begin{center}
\begin{tabular}{@{}L{3.0cm} L{10.4cm}@{}}
\toprule
\code{BEQ}, \code{BNE}   & Branch if equal / not equal. \\
\code{BLT}, \code{BGE}   & Branch if less-than / greater-or-equal (signed). \\
\code{BLTU}, \code{BGEU} & Branch if less-than / greater-or-equal (unsigned). \\
\code{BAL}               & Branch and link: PC-relative jump, save return address to a register. \\
\code{BALR}              & Branch and link to register: jump to an address held in a register, save return address. \\
\code{HLT}               & Halt. \\
\code{NOP}               & No operation. \\
\bottomrule
\end{tabular}
\end{center}

\subsection{Privilege}
\begin{center}
\begin{tabular}{@{}L{3.0cm} L{10.4cm}@{}}
\toprule
\code{ECALL}            & Trap from user mode to supervisor. \\
\code{ERET}             & Return from the trap handler to user mode. \\
\code{MFCR rd, <cr>}    & Move from a control register into a general-purpose register. \\
\code{MTCR <cr>, rs}    & Move to a control register from a general-purpose register. \\
\bottomrule
\end{tabular}
\end{center}

% ===========================================================================
\section{Instruction formats}
% ===========================================================================

Every instruction is 32 bits wide. The opcode occupies bits \code{[31:26]} in
all formats.

\subsection*{R --- register operations}
\begin{center}
\begin{bytefield}{32}
  \bitheader{0,13,17,21,25,31} \\
  \bitbox{6}{opcode} \bitbox{4}{rd} \bitbox{4}{rs1} \bitbox{4}{rs2} \bitbox{14}{unused}
\end{bytefield}
\end{center}
\noindent Used by: \code{ADD}, \code{SUB}, \code{MUL}, \code{AND}, \code{OR},
\code{XOR}, \code{SLL}, \code{SLR}, \code{SAR}, \code{MOV}, \code{NOT},
\code{BALR}. For \code{BALR}, \code{rd} is the link destination, \code{rs1} is
the branch target, and \code{rs2} is unused.

\subsection*{I --- two registers + immediate}
\begin{center}
\begin{bytefield}{32}
  \bitheader{0,17,21,25,31} \\
  \bitbox{6}{opcode} \bitbox{4}{r1} \bitbox{4}{r2} \bitbox{18}{imm}
\end{bytefield}
\end{center}
\noindent Loads: \code{r1} = destination, \code{r2} = base, \code{imm} = offset.
Stores: \code{r1} = source, \code{r2} = base, \code{imm} = offset.
Branches: \code{r1}, \code{r2} = the two operands to compare, \code{imm} = signed
PC-relative offset ($\pm$512\,KB reach).

\subsection*{U --- one register + large immediate}
\begin{center}
\begin{bytefield}{32}
  \bitheader{0,21,25,31} \\
  \bitbox{6}{opcode} \bitbox{4}{rd} \bitbox{22}{imm}
\end{bytefield}
\end{center}
\noindent \code{LDI16} / \code{LDI32L} / \code{LDI32H} use a 16-bit immediate
(upper bits unused). \code{BAL} uses the full 22-bit signed PC-relative offset
($\pm$8\,MB reach).

\subsection*{C --- control register access}
\begin{center}
\begin{bytefield}{32}
  \bitheader{0,17,21,25,31} \\
  \bitbox{6}{opcode} \bitbox{4}{rd/cr} \bitbox{4}{cr/rs} \bitbox{18}{unused}
\end{bytefield}
\end{center}
\noindent \code{MFCR}: \code{rd} = destination GPR, \code{cr} = control register
index 0--3. \code{MTCR}: \code{cr} = control register index 0--3, \code{rs} =
source GPR.

\subsection*{N --- no operands}
\begin{center}
\begin{bytefield}{32}
  \bitheader{0,25,31} \\
  \bitbox{6}{opcode} \bitbox{26}{unused}
\end{bytefield}
\end{center}
\noindent Used by \code{HLT}, \code{NOP}, \code{ECALL}, \code{ERET}.

% ===========================================================================
\section{Opcode table}
% ===========================================================================

\begin{center}
\begin{longtable}{@{}C{1.3cm} C{1.3cm} L{3.2cm} C{2.0cm}@{}}
\toprule
\textbf{Dec} & \textbf{Hex} & \textbf{Mnemonic} & \textbf{Format} \\
\midrule
\endfirsthead
\toprule
\textbf{Dec} & \textbf{Hex} & \textbf{Mnemonic} & \textbf{Format} \\
\midrule
\endhead
\bottomrule
\endfoot
0  & 0x00 & \code{HLT}    & N \\
1  & 0x01 & \code{NOP}    & N \\
2  & 0x02 & \code{MOV}    & R \\
3  & 0x03 & \code{LDI16}  & U \\
4  & 0x04 & \code{LDI32L} & U \\
5  & 0x05 & \code{LDI32H} & U \\
6  & 0x06 & \code{LDM8}   & I \\
7  & 0x07 & \code{LDM16}  & I \\
8  & 0x08 & \code{LDM32}  & I \\
9  & 0x09 & \code{STR8}   & I \\
10 & 0x0A & \code{STR16}  & I \\
11 & 0x0B & \code{STR32}  & I \\
12 & 0x0C & \code{ADD}    & R \\
13 & 0x0D & \code{SUB}    & R \\
14 & 0x0E & \code{MUL}    & R \\
15 & 0x0F & \code{AND}    & R \\
16 & 0x10 & \code{OR}     & R \\
17 & 0x11 & \code{XOR}    & R \\
18 & 0x12 & \code{NOT}    & R \\
19 & 0x13 & \code{SLL}    & R \\
20 & 0x14 & \code{SLR}    & R \\
21 & 0x15 & \code{SAR}    & R \\
22 & 0x16 & \code{BEQ}    & I \\
23 & 0x17 & \code{BNE}    & I \\
24 & 0x18 & \code{BLT}    & I \\
25 & 0x19 & \code{BGE}    & I \\
26 & 0x1A & \code{BLTU}   & I \\
27 & 0x1B & \code{BGEU}   & I \\
28 & 0x1C & \code{BAL}    & U \\
29 & 0x1D & \code{ECALL}  & N \\
30 & 0x1E & \code{ERET}   & N \\
31 & 0x1F & \code{MFCR}   & C \\
32 & 0x20 & \code{MTCR}   & C \\
33 & 0x21 & \code{BALR}   & R \\
\end{longtable}
\end{center}
\noindent Opcodes 34--63 are reserved.

% ===========================================================================
\section{Pseudo-instructions}
% ===========================================================================

Pseudo-instructions are expanded by the assembler into one or more real
instructions.

\begin{blissnote}[Clobbered registers]
Most of these use \reg{r0} as a temporary and clobber it --- do not hold a live
value in \reg{r0} across a \code{PUSH}, \code{POP}, \code{RET}, or \code{ADDI}.
\code{CALL} instead overwrites \reg{r14} (the link register) with the return
address.
\end{blissnote}

\begin{center}
\begin{tabular}{@{}L{2.6cm} L{9.2cm} C{1.3cm}@{}}
\toprule
\textbf{Pseudo} & \textbf{Expands to} & \textbf{Clob.} \\
\midrule
\code{LI rd, imm}   & \code{LDI32L rd, imm[15:0]} $\cdot$ \code{LDI32H rd, imm[31:16]} & --- \\
\code{PUSH rs}      & \code{LDI16 r0, 4} $\cdot$ \code{SUB r13, r13, r0} $\cdot$ \code{STR32 rs, [r13]} & \code{r0} \\
\code{POP rd}       & \code{LDM32 rd, [r13]} $\cdot$ \code{LDI16 r0, 4} $\cdot$ \code{ADD r13, r13, r0} & \code{r0} \\
\code{CALL label}   & \code{BAL r14, label} & \code{r14} \\
\code{RET}          & \code{BALR r0, r14} & \code{r0} \\
\code{JMP label}    & \code{BAL r0, label} & \code{r0} \\
\code{JMP rs}       & \code{BALR r0, rs} & \code{r0} \\
\code{ADDI rd, imm} & \code{LDI16 r0, imm} $\cdot$ \code{ADD rd, rd, r0} \; (imm fits 16 bits) & \code{r0} \\
\bottomrule
\end{tabular}
\end{center}

\noindent \code{CALL} is a single \code{BAL r14, label}: it sets \reg{r14} to the
return address and jumps. \code{RET} is a single \code{BALR r0, r14}: it jumps to
whatever \reg{r14} holds. There is no implicit stack traffic --- saving \reg{r14}
is the callee's job (see below).

% ===========================================================================
\section{Calling convention}
% ===========================================================================

\subsection{Register roles}

\begin{center}
\begin{tabular}{@{}L{2.4cm} L{6.8cm} L{3.4cm}@{}}
\toprule
\textbf{Register} & \textbf{Role} & \textbf{Saved by} \\
\midrule
\reg{r0}         & Assembler scratch / return value & --- (never saved) \\
\reg{r1}--\reg{r3} & Argument registers & Caller \\
\reg{r4}--\reg{r7} & Caller-saved temporaries & Caller \\
\reg{r8}--\reg{r12} & Callee-saved & Callee \\
\reg{r13}        & Stack pointer (\code{SP}) & Always preserved \\
\reg{r14}        & Link register (\code{LR}) & Caller \\
\reg{r15}        & Program counter (\code{PC}) & --- (not writable) \\
\bottomrule
\end{tabular}
\end{center}

\textbf{Caller-saved} registers (\reg{r0}--\reg{r7}, \reg{r14}) may be
overwritten freely by any function; a caller that needs them after a call must
spill them first. \textbf{Callee-saved} registers (\reg{r8}--\reg{r12}) must be
preserved: a function that touches one pushes it on entry and pops it before
returning.

\subsection{Stack}
\begin{itemize}[leftmargin=1.4em]
  \item Full descending: \code{SP} points to the last pushed value and grows
        toward lower addresses.
  \item \code{SP} is initialised to \code{0x1FFC} at kernel boot.
  \item All pushes and pops are 4-byte aligned.
\end{itemize}

\subsection{Call sequence}

\textbf{Leaf functions} --- those that make no further calls --- need no prologue
or epilogue. \reg{r14} is left untouched, so \code{RET} works directly:

\begin{lstlisting}
bar:
  ; ... body ...
  ret               ; r14 still holds the caller's return address
\end{lstlisting}

\textbf{Non-leaf functions} must save \reg{r14} before their own \code{CALL}
clobbers it, and restore it before returning. Callee-saved registers
(\reg{r8}--\reg{r12}) follow the same push-on-entry / pop-before-return pattern:

\begin{lstlisting}
foo:
  push lr           ; save caller's return address (lr = r14)
  push r8           ; save any callee-saved registers used

  ; ... body, including calls to other functions ...
  call bar

  pop r8            ; restore in reverse order
  pop lr
  ret               ; jump to the restored return address
\end{lstlisting}

% ===========================================================================
\section{Syscall ABI}
% ===========================================================================

User programs trap into the kernel with \code{ECALL}.

\begin{itemize}[leftmargin=1.4em]
  \item \reg{r0} --- syscall number on entry to the handler.
  \item \reg{r1}, \reg{r2}, \reg{r3} --- arguments.
  \item \reg{r0} --- return value, set by the handler before \code{ERET}.
\end{itemize}

\subsection{Cause codes}
\begin{center}
\begin{tabular}{@{}C{1.6cm} L{9.0cm}@{}}
\toprule
\textbf{\code{cause}} & \textbf{Meaning} \\
\midrule
0 & \code{ECALL} --- syscall from user mode. \\
\bottomrule
\end{tabular}
\end{center}

\subsection{Syscall table}
\begin{center}
\begin{tabular}{@{}C{0.9cm} L{2.2cm} L{2.6cm} L{2.6cm} L{4.0cm}@{}}
\toprule
\textbf{r0} & \textbf{Name} & \textbf{r1} & \textbf{r2} & \textbf{Returns} \\
\midrule
1 & \code{write} & buf ptr & byte count & --- \\
2 & \code{read}  & ---     & ---        & \code{r0} = byte read (blocks until available) \\
\bottomrule
\end{tabular}
\end{center}

% ===========================================================================
\section{BlissFS --- the filesystem}
% ===========================================================================

A minimal custom filesystem for the Bliss storage device.

\subsection{Parameters}
\begin{center}
\begin{tabular}{@{}L{4.2cm} L{6.6cm}@{}}
\toprule
\textbf{Property} & \textbf{Value} \\
\midrule
Block size     & 512 bytes (one sector) \\
Magic number   & \code{0xB2155F2D} \\
Max inodes     & 32 \\
Max filename   & 24 bytes (null-padded) \\
Max file size  & $8 \times 512 = 4096$ bytes \\
\bottomrule
\end{tabular}
\end{center}

\subsection{Disk layout}
\begin{center}
\begin{tabular}{@{}L{3.4cm} L{8.0cm}@{}}
\toprule
\textbf{Block} & \textbf{Contents} \\
\midrule
0      & Superblock \\
1--4   & Inode table ($32$ inodes $\times$ 64 bytes = 2048 bytes) \\
5      & Free-block bitmap \\
6+     & Data blocks \\
\bottomrule
\end{tabular}
\end{center}

\subsection{Superblock (block 0, first 20 bytes)}
\begin{center}
\begin{tabular}{@{}C{1.4cm} C{1.2cm} L{7.4cm}@{}}
\toprule
\textbf{Offset} & \textbf{Size} & \textbf{Field} \\
\midrule
0  & 4 & Magic (\code{0xB2155F2D}) \\
4  & 4 & Block size (512) \\
8  & 4 & Inode count (32) \\
12 & 4 & Data start block (6) \\
16 & 4 & Free block count \\
\bottomrule
\end{tabular}
\end{center}

\subsection{Inode (64 bytes)}
\begin{center}
\begin{tabular}{@{}C{1.4cm} C{1.2cm} L{8.4cm}@{}}
\toprule
\textbf{Offset} & \textbf{Size} & \textbf{Field} \\
\midrule
0  & 4  & File size in bytes \\
4  & 1  & Type: 0 = unused, 1 = file, 2 = directory \\
5  & 3  & Reserved \\
8  & 32 & $8 \times$ 4-byte direct block pointers (0 = unused) \\
40 & 24 & Reserved \\
\bottomrule
\end{tabular}
\end{center}
\noindent Inode 0 is reserved (null). Inode 1 is the root directory.

\subsection{Directory entry (28 bytes)}
\begin{center}
\begin{tabular}{@{}C{1.4cm} C{1.2cm} L{8.4cm}@{}}
\toprule
\textbf{Offset} & \textbf{Size} & \textbf{Field} \\
\midrule
0  & 24 & Filename, null-padded \\
24 & 4  & Inode number (0 = empty slot) \\
\bottomrule
\end{tabular}
\end{center}
\noindent 18 directory entries fit per block ($18 \times 28 = 504$ bytes; 8
bytes unused).

% ===========================================================================
\section{Worked examples}
% ===========================================================================

\subsection{Hello, world (bare metal, no OS)}

Writes a string to the serial TX port one byte at a time.

\lstinputlisting{../examples/hello.asm}

\subsection{A user-mode program}

Loaded by the kernel at \code{0x2000}, prints via the \code{write} syscall.

\lstinputlisting{../os/user.asm}

\end{document}
